\section{Introduction}\label{sec:introduction}

Dispersion and diversity problems arise in the placement of undesirable facilities, workforce management and social media, among others \citep{peiro2021}; our application-inspired MDP instances mimic sensor, facility, portfolio and team selection. The classical formulations optimise one measure of dispersion: MaxSum (maximise the sum of pairwise distances, the maximum diversity problem), or MaxMin (maximise the minimum pairwise distance, the p-dispersion family). The two measures lead to solutions of very different structure --- the MaxSum optimum selects points close to the borders of the space, whereas the MaxMin optimum spreads them evenly over it, central region included \citep{parreno2021} --- and practitioners often care about both. Two constrained variants add realism: the capacitated dispersion problem (CDP) requires the selected elements to supply a minimum aggregate capacity, and the generalised dispersion problem (GDP) additionally bounds the total opening cost. The bi-objective version of the unconstrained problem, pairing MaxSum with MaxMin, is the \emph{bi-objective diversity problem} of \citet{colmenar2018}, which they attacked with NSGA-II, SPEA2, GRASP, iterated greedy, tabu search and VNS on the 145 GKD instances of MDPLIB (n up to 500, m up to 50), each method tuned with irace \citep{lopezibanez2016}; to our knowledge the literature on CDP and GDP, in contrast, is single-objective (MaxMin), and the one multi-objective CDP study we are aware of pairs MaxMin with a colour-symmetry penalty rather than with MaxSum and uses an $\varepsilon$-constraint method and a biased-randomised constructive heuristic \citep{gomez2026}.

Decomposition-based evolutionary algorithms such as MOEA/D \citep{zhang2007} are natural candidates, and recent work equips them with \emph{adaptive operator selection} (AOS): a bandit chooses, at every step, which variation operator to apply, rewarded by the improvement it produces (DaCosta et al., 2008, dynamic bandits; Li et al., 2014, MOEA/D-FRRMAB, with a sliding window of fitness-improvement rates, rank-based credit assignment and a decaying selection rule; Sun and Li, 2020, dynamic Thompson sampling; Geng et al., 2023, MOEA/D-CTAOS, with a classification tree that predicts the outcome of each operator). The premise is that operators differ in usefulness and that this can be learned online. This paper tests that premise in a controlled setting and asks four questions.

\begin{itemize}
\item \textbf{RQ1.} Does a MOEA/D framework with a pool of eight operators outperform classical multi-objective heuristics (NSGA-II, GRASP, VNS, tabu search) on bi-objective MDP, CDP and GDP --- also when those heuristics are tuned and given the same wall-clock time?
\item \textbf{RQ2.} Does adaptive selection add value over \emph{uniform} selection from the same pool, and how does that depend on instance size, constraint tightness and budget?
\item \textbf{RQ3.} How much do the scaling of objectives and the coverage of the Pareto-front ends matter?
\item \textbf{RQ4.} Which parts of the operator pool matter, and is the pool robust to pruning?
\end{itemize}

\textbf{Contributions.}

\begin{itemize}
\item \emph{Protocol.} A replicated, indicator-robust evaluation of the bi-objective MDP, CDP and GDP on 443 benchmark instances, against frozen reference fronts and with three quality indicators, validated against the public single-objective benchmarks. It includes equal-time comparisons with baselines --- NSGA-II, SMS-EMOA, GRASP, and a tabu search and a VNS in the style of Colmenar et al. --- that are tuned by the same protocol as the proposed MOEA/D.
\item \emph{Where MOEA/D stands.} The answer depends on the problem and on how the baselines are treated. MOEA/D is clearly ahead on the GDP and, at equal time, on the CDP, where five times the time lets tuned GRASP, VNS and tabu search catch up; it is behind a tabu search, a VNS and a tuned GRASP on the unconstrained MDP.
\item \emph{Objective scaling.} Normalising the objectives inside the Tchebycheff function is the single most important design choice, with an effect that grows with instance size.
\item \emph{Adaptive operator selection.} A controlled, replicated study with three bandit families, the published dynamic Thompson sampling and FRRMAB, several reward signals, hyper-parameters, budgets and contexts. Uniform selection from the pool is as good as any selector on instances up to about 200 elements and on the CDP, and a sizeable family of selectors is significantly worse. The HV-reward selectors do win on large instances and on instances with a tight cost budget, with an advantage that depends on the generation budget: a map of when adaptive selection pays off, not a blanket verdict.
\item \emph{The pool.} No operator suffices alone and MaxMinSwap cannot be removed, but BiObjBestSwap, BestSwap1$\times$1 and GraspGreedy can be (singly on both problems, together on the CDP) without loss, and the uniform pool also tolerates three deliberately useless operators added to it.
\item \emph{Covering the ends of the front.} A short MaxMin intensification (a tabu search followed by a solution-based tabu search) brings the MaxMin end of the front to within about 1\% of the best published single-objective values on four of five CDP families and above the published GDP values, and improves the hypervolume; the hardest family (CDP MDG-b) remains at 0.85 of the best published value.
\item \emph{A hidden dimension.} In the CDP and GDP the size of the selected set is a hidden dimension of the bi-objective problem, which all methods leave at the size induced by the capacity constraint; we quantify what this leaves unexplored.
\item \emph{Open data.} Reference fronts, about 100,000 run files and the analysis scripts that regenerate every table and figure.
\end{itemize}
